\documentclass[twocolumn,amsmath,amssymb,aps,prl,superscriptaddress]{revtex4-2}

\usepackage{amsmath,amssymb,amsthm}
\usepackage{graphicx}
\usepackage{hyperref}
\usepackage{physics}
\usepackage{bm}
\usepackage{bbm}
\usepackage{mathtools}
\usepackage{xcolor}
\usepackage{booktabs}

% Theorem environments
\newtheorem{theorem}{Theorem}[section]
\newtheorem{lemma}[theorem]{Lemma}
\newtheorem{proposition}[theorem]{Proposition}
\newtheorem{corollary}[theorem]{Corollary}
\newtheorem{definition}[theorem]{Definition}
\newtheorem{axiom}{Axiom}
\newtheorem{principle}{Principle}

% Custom commands
\newcommand{\Sspace}{\mathcal{S}}
\newcommand{\Sk}{S_k}
\newcommand{\St}{S_t}
\newcommand{\Se}{S_e}
\newcommand{\Mstate}{\mathcal{M}}
\newcommand{\Ccat}{\mathcal{C}}
\newcommand{\Pdecay}{P_{\text{decay}}}
\newcommand{\Tdecay}{T_{\text{decay}}}
\newcommand{\Hfield}{H}
\newcommand{\Gstate}{\Gamma}
\newcommand{\traj}{\gamma}

\begin{document}

\title{Psychon Circuit Mechanics: \\
S-Entropy Coordinates, Trajectory States, and Consciousness Geometry}

\author{Blickrichtung Collaboration}
\affiliation{St Stellas Laboratory for Oscillatory Systems}

\date{\today}

\begin{abstract}
We develop a circuit-theoretic framework for mental activity based on the \emph{psychon}: a trajectory-terminus pair $\Mstate = (\traj, \Gstate_f)$ in S-entropy space. The S-entropy coordinates $(\Sk, \St, \Se) \in [0,1]^3$ provide a universal basis for categorical states, where $\Sk$ measures information deficit, $\St$ temporal position, and $\Se$ categorical entropy. Each psychon operates simultaneously as resistor (in $\Sk$), capacitor (in $\St$), and inductor (in $\Se$), with actual mode selected by S-entropy minimization. We prove that \emph{consciousness} $\Ccat(t) = \Pdecay(t) \cap \Tdecay(t)$ is the geometric intersection of perception and thought decay curves, with magnitude $|\Ccat| = \min(\Pdecay, \Tdecay)$. From this framework emerge: (i) mental state non-identity for identical thoughts via different trajectories, (ii) memory as accumulated emotional field $M = \int \dot{\Hfield}^+ dt$, (iii) dreams as unconstrained thought evolution ($\Pdecay = 0$), and (iv) pharmacological effects as trajectory modification. The framework provides a geometric rather than computational theory of mental phenomena.
\end{abstract}

\maketitle

%==============================================================================
\section{Introduction}
%==============================================================================

The relationship between neural activity, information processing, and conscious experience remains one of science's central problems. Computational approaches treat the brain as information processor; correlational approaches map neural activity to behavior; neither provides a geometric understanding of mental states.

We present a framework grounded in three principles: (i) mental states are trajectory-terminus pairs in a coordinate space, (ii) these coordinates form a tri-dimensional S-entropy system with circuit equivalents, and (iii) consciousness is the geometric intersection of perception and thought decay curves. The framework treats mental phenomena as geometric structure rather than computational process.

The consequences are substantial. The same thought content constitutes different mental states when reached via different trajectories. Memory is the time-integral of emotional change. Dreams are unconstrained thought evolution without perceptual grounding. Drugs modify mental states by altering trajectories, not termini.

This paper establishes the mathematical foundations. Section~\ref{sec:sentropy} introduces S-entropy coordinates. Section~\ref{sec:psychon} defines the psychon as trajectory-terminus pair. Section~\ref{sec:circuits} develops circuit equivalence. Section~\ref{sec:consciousness} proves consciousness as intersection. Section~\ref{sec:consequences} derives mental phenomena.

%==============================================================================
\section{S-Entropy Coordinate System}
\label{sec:sentropy}
%==============================================================================

\subsection{The Tri-Dimensional Basis}

\begin{definition}[S-Knowledge]
\label{def:sk}
The knowledge entropy $\Sk$ measures information deficit:
\begin{equation}
\boxed{\Sk = -\log P(\text{equivalence class})}
\end{equation}
where $P$ is the probability of the categorical equivalence class. Higher $\Sk$ indicates greater uncertainty about categorical state.
\end{definition}

\begin{definition}[S-Time]
\label{def:st}
The temporal entropy $\St$ measures position in categorical completion:
\begin{equation}
\boxed{\St = \frac{\text{steps to completion}}{\text{total possible steps}} \in [0, 1]}
\end{equation}
$\St = 0$ indicates completion; $\St = 1$ indicates initiation.
\end{definition}

\begin{definition}[S-Entropy]
\label{def:se}
The categorical entropy $\Se$ measures distribution over states:
\begin{equation}
\boxed{\Se = -\sum_i p_i \log p_i}
\end{equation}
the Shannon entropy of the state distribution.
\end{definition}

\begin{theorem}[Coordinate Sufficiency]
\label{thm:sufficient}
The coordinates $(\Sk, \St, \Se)$ form a sufficient basis for categorical state specification in bounded systems.
\end{theorem}

\begin{proof}
By the Bounded Phase Space Law, all physical systems occupy finite phase space. Any finite partition admits characterization by: (i) which cell is occupied ($\Sk$), (ii) completion progress ($\St$), and (iii) residual uncertainty ($\Se$). Additional dimensions provide refinement but not fundamentally new information.
\end{proof}

\subsection{S-Entropy Space}

\begin{definition}[S-Entropy Space]
The S-entropy space is the unit cube:
\begin{equation}
\Sspace_N = [0, 1]^3 \ni (\Sk, \St, \Se)
\end{equation}
with natural metric:
\begin{equation}
d_S(\mathbf{s}_1, \mathbf{s}_2) = \|\mathbf{s}_1 - \mathbf{s}_2\|
\end{equation}
\end{definition}

\begin{theorem}[Cross-Domain Equivalence]
\label{thm:equiv}
Systems A and B are informationally equivalent if:
\begin{equation}
\|(\Sk^A, \St^A, \Se^A) - (\Sk^B, \St^B, \Se^B)\| < \varepsilon
\end{equation}
for threshold $\varepsilon \approx 0.1$.
\end{theorem}

This equivalence enables solution transfer between disparate physical domains occupying equivalent S-entropy coordinates.

\subsection{Extended Coordinates}

For refined specification, additional coordinates are available:

\begin{definition}[Extended S-Coordinates]
\begin{align}
S_{\text{packing}} &: \text{geometric configuration} \\
S_{\text{hydrophobic}} &: \text{energy landscape position}
\end{align}
The full coordinate vector is $(\Sk, \St, \Se, S_{\text{packing}}, S_{\text{hydrophobic}}) \in [0,1]^5$.
\end{definition}

The primary tri-dimensional coordinates suffice for most applications; extended coordinates provide domain-specific refinement.

%==============================================================================
\section{The Psychon: Trajectory-Terminus Pairs}
\label{sec:psychon}
%==============================================================================

\subsection{Definition}

\begin{definition}[Psychon]
\label{def:psychon}
A psychon is a trajectory-terminus pair in S-entropy space:
\begin{equation}
\boxed{\Mstate = (\traj, \Gstate_f)}
\end{equation}
where:
\begin{itemize}
\item $\traj: [0, T] \to \Sspace_N$ is the trajectory through S-entropy space
\item $\Gstate_f = \traj(T)$ is the terminus (final) state
\end{itemize}
\end{definition}

The psychon is the fundamental unit of mental activity. Crucially, it includes both the endpoint and the path taken to reach it.

\subsection{Why Trajectory Matters}

\begin{theorem}[Mental State Non-Identity]
\label{thm:nonidentity}
Different trajectories to identical termini constitute different mental states:
\begin{equation}
\boxed{\traj_1 \neq \traj_2 \implies \Mstate_1 \neq \Mstate_2 \quad \text{even if } \Gstate_{f,1} = \Gstate_{f,2}}
\end{equation}
\end{theorem}

\begin{proof}
The mental state is defined as the pair $\Mstate = (\traj, \Gstate_f)$. If $\traj_1 \neq \traj_2$, then $(\traj_1, \Gstate_f) \neq (\traj_2, \Gstate_f)$ as ordered pairs, regardless of terminus equality.
\end{proof}

\begin{corollary}[Context Dependence]
The same thought content in different contexts constitutes different mental states.
\end{corollary}

Consider the thought ``cup on chair'':
\begin{align}
\Mstate_{\text{party}} &= (\traj_{\text{festive}}, \Gstate_{\text{cup-chair}}) \\
\Mstate_{\text{exam}} &= (\traj_{\text{academic}}, \Gstate_{\text{cup-chair}}) \\
\Mstate_{\text{search}} &= (\traj_{\text{goal-directed}}, \Gstate_{\text{cup-chair}}) \\
\Mstate_{\text{memory}} &= (\traj_{\text{nostalgic}}, \Gstate_{\text{cup-chair}})
\end{align}

Four distinct mental states share identical thought content but differ in trajectory---hence in full mental state specification.

\subsection{Trajectory Context}

\begin{definition}[Trajectory Context]
The context carried by a trajectory is:
\begin{equation}
\text{Context}(\traj) = \int_0^T \frac{d\traj}{dt} \otimes \Hfield(t) \, dt
\end{equation}
where $\Hfield(t)$ is the emotional field and $\otimes$ denotes outer product.
\end{definition}

The context integral captures how emotional state modulates trajectory evolution. Identical paths traversed under different emotional conditions carry different contexts.

\subsection{Ternary Address Structure}

\begin{theorem}[Trajectory Address]
\label{thm:address}
Each trajectory corresponds to a ternary address:
\begin{equation}
\traj \leftrightarrow (a_1, a_2, \ldots, a_k) \quad \text{where } a_i \in \{0, 1, 2\}
\end{equation}
The address IS the path (trajectory-position identity).
\end{theorem}

\begin{corollary}[Navigation Efficiency]
Navigation through mental state space has complexity:
\begin{equation}
\text{Complexity} = O(\log_3 n)
\end{equation}
where $n$ is the number of distinct states. This is 37\% more efficient than binary addressing.
\end{corollary}

%==============================================================================
\section{Circuit Equivalence}
\label{sec:circuits}
%==============================================================================

\subsection{The R-C-L Trichotomy}

\begin{theorem}[Simultaneous Operation]
\label{thm:rcl}
Every psychon operates simultaneously as:
\begin{align}
\text{Resistor} &\quad \text{in } \Sk \text{ dimension} \\
\text{Capacitor} &\quad \text{in } \St \text{ dimension} \\
\text{Inductor} &\quad \text{in } \Se \text{ dimension}
\end{align}
Actual behavior is selected by S-entropy minimization.
\end{theorem}

The circuit analogy provides physical intuition:
\begin{itemize}
\item \textbf{Resistive mode}: Dissipates information deficit, reduces $\Sk$
\item \textbf{Capacitive mode}: Stores temporal position, governs $\St$ evolution
\item \textbf{Inductive mode}: Resists entropy change, stabilizes $\Se$
\end{itemize}

\subsection{Tri-Dimensional Parameters}

\begin{definition}[Circuit Parameters]
The circuit parameters are related by:
\begin{align}
R &= f(\Sk) & \text{(resistance)} \\
C &= g(\St) & \text{(capacitance)} \\
L &= h(\Se) & \text{(inductance)}
\end{align}
with constraint:
\begin{equation}
C = \frac{\tau}{\pi R}, \quad L = \frac{\pi R}{\tau}
\end{equation}
where $\tau$ is the characteristic time constant.
\end{definition}

\subsection{S-Entropy Minimization}

\begin{definition}[Mode Selection]
The weighted S-entropy cost is:
\begin{equation}
\text{Cost} = \alpha \Sk + \beta \St + \gamma \Se
\end{equation}
where $(\alpha, \beta, \gamma)$ are context-dependent weights.
\end{definition}

\begin{theorem}[Optimal Mode]
The psychon selects the mode minimizing weighted S-entropy cost:
\begin{equation}
\text{Mode}^* = \argmin_{\text{mode} \in \{R, C, L\}} \text{Cost}_{\text{mode}}
\end{equation}
\end{theorem}

In practice:
\begin{itemize}
\item High $\Sk$ (information deficit) $\to$ Resistive mode
\item High $\St$ (temporal urgency) $\to$ Capacitive mode
\item High $\Se$ (categorical diversity) $\to$ Inductive mode
\end{itemize}

\subsection{Impedance}

For frequency $\omega$, the complex impedance in each mode:
\begin{equation}
Z = \begin{cases}
R & \text{resistive} \\
-j/(\omega C) & \text{capacitive} \\
j\omega L & \text{inductive}
\end{cases}
\end{equation}

\begin{table}[h]
\centering
\caption{Measured Circuit Parameters}
\begin{tabular}{lcc}
\toprule
Parameter & Value & Units \\
\midrule
Hole mobility & 0.0123 & cm$^2$/(V$\cdot$s) \\
On/off ratio & 42.1 & -- \\
Built-in potential & 0.615 & eV \\
Hole concentration & $2.80 \times 10^{12}$ & cm$^{-3}$ \\
\bottomrule
\end{tabular}
\end{table}

%==============================================================================
\section{Consciousness as Geometric Intersection}
\label{sec:consciousness}
%==============================================================================

\subsection{Decay Curves}

\begin{definition}[Perception Decay]
The perception decay curve $\Pdecay(t)$ measures perceptual input integration:
\begin{equation}
\frac{d\Pdecay}{dt} = -\frac{\Pdecay}{\tau_P} + I_{\text{sensory}}(t)
\end{equation}
where $\tau_P \approx 50$ ms is the perceptual integration time and $I_{\text{sensory}}$ is sensory input rate.
\end{definition}

\begin{definition}[Thought Decay]
The thought decay curve $\Tdecay(t)$ measures cognitive process evolution:
\begin{equation}
\frac{d\Tdecay}{dt} = -\frac{\Tdecay}{\tau_T} + T^{\infty}_{\text{decay}} + \Hfield(t)
\end{equation}
where $\tau_T \approx 100$ ms is thought decay time, $T^{\infty}_{\text{decay}}$ is background activity, and $\Hfield$ is the emotional field.
\end{definition}

\subsection{The Intersection Theorem}

\begin{theorem}[Consciousness Intersection]
\label{thm:consciousness}
Consciousness $\Ccat$ is the geometric intersection of perception and thought decay curves:
\begin{equation}
\boxed{\Ccat(t) = \Pdecay(t) \cap \Tdecay(t)}
\end{equation}
\end{theorem}

The intersection has precise meaning: consciousness occurs when and only when both perception and thought are simultaneously active above threshold.

\begin{theorem}[Consciousness Magnitude]
\label{thm:magnitude}
The magnitude of consciousness at time $t$ is:
\begin{equation}
\boxed{|\Ccat(t)| = \min(\Pdecay(t), \Tdecay(t))}
\end{equation}
Consciousness is limited by the weaker of perception or thought.
\end{theorem}

\begin{proof}
The intersection of two sets has measure bounded by the smaller set. For decay curves as measure functions, $|\Ccat| = \min(\Pdecay, \Tdecay)$ at each point.
\end{proof}

\subsection{Consciousness Window}

\begin{definition}[Consciousness Window]
The consciousness window is the temporal region where both curves exceed threshold $\theta$:
\begin{equation}
W_C = \{t : \Pdecay(t) > \theta \land \Tdecay(t) > \theta\}
\end{equation}
\end{definition}

\begin{theorem}[Window Duration]
\label{thm:window}
The consciousness window has characteristic duration:
\begin{equation}
\Delta t_C = \frac{\tau_P \tau_T}{\tau_P + \tau_T} \ln\left(\frac{\Pdecay(0) \Tdecay(0)}{\theta^2}\right)
\end{equation}
For typical values, $\Delta t_C \sim 100$--$500$ ms.
\end{theorem}

\begin{proof}
For exponentially decaying curves with initial values $\Pdecay(0)$, $\Tdecay(0)$ and time constants $\tau_P$, $\tau_T$, the window closes when the first curve crosses $\theta$. The effective time constant for the intersection is the harmonic mean $\tau_{\text{eff}} = \tau_P \tau_T / (\tau_P + \tau_T)$.
\end{proof}

\subsection{Non-Intersection States}

\begin{proposition}[State Classification]
Mental states are classified by intersection status:
\begin{center}
\begin{tabular}{ccc}
\toprule
$\Pdecay$ & $\Tdecay$ & State \\
\midrule
$> 0$ & $> 0$ & Conscious \\
$= 0$ & $> 0$ & Dreaming \\
$> 0$ & $= 0$ & Reflexive \\
$= 0$ & $= 0$ & Unconscious \\
\bottomrule
\end{tabular}
\end{center}
\end{proposition}

Dreams represent thought evolution without perceptual constraint. Reflexes represent perceptual response without cognitive elaboration. Full consciousness requires intersection.

\subsection{Self-Reference}

\begin{corollary}[Self-Consciousness]
Self-consciousness is self-referential intersection:
\begin{equation}
\Ccat_{\text{self}} = \Pdecay(\text{self-perception}) \cap \Tdecay(\text{self-thought})
\end{equation}
\end{corollary}

When perception includes the body and thought includes self-reference, the intersection creates awareness of awareness---the self-referential loop characteristic of human consciousness.

%==============================================================================
\section{Emergent Mental Phenomena}
\label{sec:consequences}
%==============================================================================

\subsection{Memory}

\begin{definition}[Memory]
\label{def:memory}
Memory is the time-integral of positive emotional field change:
\begin{equation}
\boxed{M = \int_0^t \dot{\Hfield}^+ \, dt'}
\end{equation}
where $\dot{\Hfield}^+ = \max(d\Hfield/dt, 0)$.
\end{definition}

\begin{theorem}[Memory Encoding]
Memory formation requires emotional field change. Static emotional states do not encode memory.
\end{theorem}

\begin{proof}
If $d\Hfield/dt = 0$, then $\dot{\Hfield}^+ = 0$, and $M = 0$ over any interval. Memory requires $\dot{\Hfield}^+ > 0$.
\end{proof}

\begin{corollary}[Emotional Memory Enhancement]
Emotionally significant events produce larger $|\dot{\Hfield}|$, hence stronger memory encoding.
\end{corollary}

\subsection{Dreams}

\begin{theorem}[Dream Dynamics]
\label{thm:dreams}
Dreams are characterized by:
\begin{equation}
\Pdecay = 0, \quad \Tdecay > 0
\end{equation}
Thought evolves without perceptual constraint.
\end{theorem}

\begin{corollary}[Dream Phenomenology]
\begin{enumerate}
\item \textbf{Unconstrained association}: Without $\Pdecay$, trajectories are not reality-tested
\item \textbf{Emotional intensity}: $\Hfield$ directly drives $\Tdecay$ without perceptual modulation
\item \textbf{Time distortion}: $\St$ evolves without external temporal reference
\end{enumerate}
\end{corollary}

\subsection{Pharmacological Effects}

\begin{theorem}[Drugs as Trajectory Modifiers]
\label{thm:drugs}
Pharmacological agents modify mental states by altering trajectories:
\begin{equation}
\mathcal{P}_{\text{drug}}: \traj \to \traj'
\end{equation}
The terminus may remain unchanged while the mental state differs.
\end{theorem}

\begin{corollary}[Anxiolytic Effect]
An anxiolytic modifies trajectory without changing thought content:
\begin{align}
\traj_{\text{anxious}} &\to \traj_{\text{calm}} \\
\Gstate_f &= \Gstate_f \quad \text{(same thought)} \\
\Mstate_{\text{anxious}} &\neq \Mstate_{\text{calm}} \quad \text{(different mental states)}
\end{align}
\end{corollary}

\begin{table}[h]
\centering
\caption{Pharmacological Effects on Decay Parameters}
\begin{tabular}{lccc}
\toprule
Drug Class & $\tau_P$ & $\tau_T$ & Effect \\
\midrule
Anxiolytics & $\uparrow$ & $\uparrow$ & Wider window \\
Stimulants & $\downarrow$ & $\downarrow$ & Narrower, intense \\
Psychedelics & Variable & $\uparrow\uparrow$ & Altered geometry \\
Anesthetics & $\to 0$ & $\to 0$ & No intersection \\
\bottomrule
\end{tabular}
\end{table}

\subsection{Penultimate State Principle}

\begin{principle}[Penultimate Determination]
To identify a mental state, determine the penultimate state from which it was reached:
\begin{equation}
\Mstate = (\traj, \Gstate_f) \iff \Gstate_{f-1} \xrightarrow{\delta\traj} \Gstate_f
\end{equation}
\end{principle}

This inverts the standard paradigm: instead of asking ``where will this go?'' (forward simulation), we ask ``where did this come from?'' (backward constraint satisfaction).

\begin{theorem}[Backward Uniqueness]
Given $\Gstate_f$ and the constraint structure, the penultimate state $\Gstate_{f-1}$ is unique:
\begin{equation}
\Gstate_f + \text{constraints} \implies \Gstate_{f-1} \quad \text{(unique)}
\end{equation}
\end{theorem}

\subsection{Mental State Equivalence}

\begin{definition}[Trajectory Equivalence]
Two trajectories are equivalent if they reach the same terminus with equivalent contexts:
\begin{equation}
\traj_1 \sim \traj_2 \iff \Gstate_{f,1} = \Gstate_{f,2} \land |\text{Context}(\traj_1) - \text{Context}(\traj_2)| < \varepsilon
\end{equation}
\end{definition}

\begin{theorem}[Equivalence Classes]
Mental states partition into equivalence classes based on:
\begin{enumerate}
\item Terminus identity: Same $\Gstate_f$
\item Context similarity: Similar trajectory integrals
\item Emotional proximity: Similar $\Hfield$ profiles
\end{enumerate}
\end{theorem}

%==============================================================================
\section{Discussion}
%==============================================================================

The framework presented here establishes mental states as trajectory-terminus pairs in S-entropy space, with consciousness as the geometric intersection of perception and thought decay curves.

Several features merit emphasis:

\textbf{Geometric Theory.} Mental phenomena are geometric structures, not computational processes. Consciousness is intersection; memory is integral; dreams are unconstrained evolution. No information processing metaphor is required.

\textbf{Trajectory Primacy.} The path matters as much as the destination. This explains context-dependence, state-dependent memory, and why the same thought ``feels different'' in different circumstances.

\textbf{Quantitative Predictions.} The framework predicts: (i) consciousness window duration $\sim 100$--$500$ ms, (ii) memory encoding proportional to emotional change rate, (iii) drug effects via decay parameter modification.

\textbf{Circuit Correspondence.} The R-C-L trichotomy provides physical intuition and potential implementation pathways for artificial systems.

%==============================================================================
\section{Conclusions}
%==============================================================================

We have established psychon circuit mechanics:

\begin{enumerate}
\item \textbf{S-Entropy Coordinates}: $(\Sk, \St, \Se) \in [0,1]^3$ provide a universal basis for categorical states.

\item \textbf{Psychon}: The fundamental unit of mental activity is the trajectory-terminus pair $\Mstate = (\traj, \Gstate_f)$.

\item \textbf{Circuit Equivalence}: Psychons operate as R/C/L circuits in the three S-dimensions, with mode selected by S-entropy minimization.

\item \textbf{Consciousness}: $\Ccat = \Pdecay \cap \Tdecay$ is the geometric intersection of decay curves, with magnitude $|\Ccat| = \min(\Pdecay, \Tdecay)$.

\item \textbf{Consequences}: Memory, dreams, and pharmacological effects emerge as natural consequences of the geometry.
\end{enumerate}

The framework provides a mathematical treatment of mental phenomena without recourse to computational metaphor, grounding psychology in the geometry of categorical state spaces.

\begin{acknowledgments}
This work was developed within the St Stellas Laboratory theoretical program on consciousness and categorical mechanics.
\end{acknowledgments}

\bibliography{references}

\end{document}